% Page 079

\seite{79}

\abschnitt{1926}
\pstart
Voll Sorgen und Kummer traten wir\ob
mit dem verflossenen in das neue Jahr. Ein Blick\ob
in die Welt, besonders nach Deutschland, läßt uns\ob
das ganze Elend vor Augen. Kein Geld, kein Brot\ohb
Verdienste und Not, Verzweiflung und Entbehrung\ob
Geschäftskrisen, Konkurs über Konkurs, Bank\ohb
rott über Bankrott! Ein Millionenheer von Arbeits\ohb
losen! Hohe Preise! Keine Aussicht auf Besserung!\ob
Himmel, hilf, sonst fallen die Sorgen über uns! —

15.\ Januar.  Wie im vorigen, so ist auch diesen Winter eine\ob
Fortbildungsschule hier eröffnet worden. Die Teil\ohb
nehmerzahl beträgt 13, schulentlassene Burschen, als\ob
Arbeitsleistungen des Tagesunterrichts finden je\ohb
weils Montags, Mittwochs und Samstags von\ob
7 bis 9½ Uhr abends statt.

3.\ März      In stiller Trauer gedachten wir unserer Volksgenossen\ob
am 28.\ Februar, dem allgemeinen Volkstrauer\ohb
tag, jener 7 im Weltkriege gefallenen Brüder. Die\ob
Fahnen wurden halbmast geflaggt. Das Glöcklein der\ob
Kapelle verkündete seine Teilnahme an der Volks\ohb
trauer von 1 – 1½ Uhr mittags. Die Gedenktafel war\ob
würdig geziert. Des Abends Rosenkranz ward den\ob
Gefallenen gewidmet.

Infolge der milden Witterung mußte auf Wunsch\ob
der älteren Teilnehmer der Fortbildungsschule / Winterschule\ob
erst am 1.\ März geschlossen werden, weil die Feld\ohb
arbeiten bereits im Gange sind. Die so\ohb
genannte Teilnehmerzahl beträgt 100\unsicher.

g/. 16.3.26.

\pend
\begin{verse}
Gaubach. L.
\end{verse}
\pstart
\pend
